Drones show great potential for expanding the reach of data collection into remote, 
hazardous, or otherwise difficult to reach areas. However, their use is limited by their high power 
consumption and resulting short mission time. One method of extending a drone's mission time is through
perching behavior, in which a drone is capable of attaching to a stationary object or surface while cutting 
off power to its motors. Various methods of implementing this behavior have been researched, though many 
are mechanically complex, require complex controls, or are only suitable for very small drones. This thesis 
proposes a perch detection mechanism and grasping device that is low cost, simple to fabricate, and can be
implemented with free, open-sourced software libraries. The perch detection method is demonstrated to function 
under ideal conditions, with further development needed to increase its reliability across different backgrounds.